%% kapa: exam=1 round=28 year=2017 period=1 form=A subject=economics subject_ko=경제학원론 seq=39 number=79
%% source: Q-Net 감정평가사 기출문제 (cst00309) artlSeq=5199603 file=extracted/1008/5199603/2017년 제28회 감정평가사 1차시험 문제지/1교시 A형.pdf
%% conversion: kordoc 4.24.51 -> tools/gen1.py (auto); review=proofread-sonnet; flags=none
\begin{question}{79}
리카디언 등가정리(Ricardian equivalence theorem)가 성립할 경우 옳은 설명을 모두 고른 것은?
\begin{framedtext}
ㄱ. 현재소비는 기대되는 미래소득과 현재소득을 모두 포함한 평생소득(lifetime income)에 의존한다. \par
ㄴ. 소비자는 현재 차입제약 상태에 있다. \par
ㄷ. 다른 조건이 일정할 때, 공채발행을 통한 조세삭감은 소비에 영향을 줄 수 없다. \par
ㄹ. 정부지출 확대정책은 어떠한 경우에도 경제에 영향을 줄 수 없다.
\end{framedtext}
\begin{choices}
  \item ㄱ, ㄷ
  \item ㄱ, ㄹ
  \item ㄴ, ㄷ
  \item ㄱ, ㄷ, ㄹ
  \item ㄴ, ㄷ, ㄹ
\end{choices}
\end{question}
\answer{1}
